% SPDX-License-Identifier: Apache-2.0
% covers: Arbiter2, Arbiter3, Arbiter4, Arbiter5, Arbiter6, Arbiter7, Arbiter8
\datasheet{Arbiter}{AXI4 arbiters: two to eight hosts merged onto one peripheral link, answered by identifier.}

\dsfacts{\code{lib/parts/src/bus/arbiter.rs}; generator \code{arbiter\_text} in \code{lib/macros/lib.rs}}%
{\code{txhdl\_parts::bus::arbiter::Arbiter2} to \code{Arbiter8}}%
{\code{//docs:axi}}

\subsection*{Function}

An arbiter puts several hosts on one peripheral link. It is the
mirror of the \code{Router}: the router decides on the address and
fans out, the arbiter decides whose turn it is and merges. It takes
\code{aw}\textit{i}, \code{ar}\textit{i} and \code{w}\textit{i} from
each host \textit{i} and gives the peripheral side one \code{aw}, one
\code{ar} and one \code{w}; it takes that side's \code{b} and
\code{r} and gives each host its own.

Two things follow from AXI4 rather than from any choice made here.

An answer must reach the host that asked, and the peripheral does not
know there is more than one. The only thing it carries back is the
identifier from the address phase, so the arbiter writes the host's
port number above the host's own identifier on the way out and takes
it off again on the way back. The peripheral side's identifier is
therefore wider than a host's.

The write data channel carries no identifier at all: a beat belongs
to the oldest write address phase that has not finished. So the
arbiter holds that channel for the host whose write address phase it
granted until that burst's last beat, and a second host's address
phase waits. Reads need no such lock.

\code{arbiter!(ArbiterN, N)} writes one unit per count of hosts,
because the lowering reads a body and not a loop over ports.
\code{arbiter.rs} writes \code{Arbiter2} to \code{Arbiter8}; the macro
refuses a count outside 2 to 8.

\subsection*{Parameters}

\begin{center}\footnotesize
\begin{tabular}{@{}lp{0.68\textwidth}@{}}
\toprule
Parameter & Meaning \\
\midrule
\code{A} & The address width, in bits. \\
\code{D} & The data width, in bits. \\
\code{S} & The write strobe width, one bit per byte lane, so \code{D / 8}. \\
\code{I} & The identifier width on each host, in bits. \\
\code{J} & The identifier width on the peripheral side. It must be at
least \code{I} plus the bits the port number needs, which is one bit
for two hosts and three for five to eight. Nothing checks it; too
small a \code{J} loses the top of the tag. \\
\code{FIXED} & Zero for round robin, anything else for fixed
priority. \\
\bottomrule
\end{tabular}
\end{center}

\subsection*{The identifier}

A host's address phase goes out carrying
\code{(port << I) | id}, where \code{id} is the identifier that host's
own tracker allocated. An answer's identifier is read the other way:
\code{id >> I} names the host, and its low \code{I} bits are given
back to that host unchanged.

So a host's identifiers stay its own, and what the peripheral sees is
a wider space with the hosts side by side in it. The tables below are
generated for \code{Arbiter2<16, 32, 4, 2, 5, 0>}, the example's, whose
hosts carry two bits and whose peripheral side carries five.

\dstables{Arbiter}

\subsection*{Behaviour}

\begin{itemize}
\item \textbf{Read address.} Each cycle the hosts offering on
\code{ar} are considered, and one is granted if the peripheral side
has room. Several reads may be outstanding at once, from different
hosts, since each carries its identifier.
\item \textbf{Write address.} The same, except that no host is granted
while \code{wbusy} is set. Taking a phase sets \code{wbusy} and
records the host in \code{wsel}, a bit each.
\item \textbf{Write beats.} While \code{wbusy} is set, beats pass from
the host in \code{wsel} and from no other. The beat marked \code{last}
clears \code{wbusy}, and the next write address phase may then be
granted, to any host.
\item \textbf{Answers.} A beat on \code{b} or \code{r} is given to the
host its identifier's top bits name, when that host has room, with the
tag taken off. A host that is not ready holds that beat, and the
peripheral side is not taken from until it is.
\item \textbf{The turn.} With \code{FIXED} zero, \code{rturn} and
\code{wturn} hold where to start looking, and each moves past the host
that won, so a host that keeps asking cannot shut another out. With
\code{FIXED} not zero every host is eligible always, so the lowest
numbered one that is offering wins.
\end{itemize}

\subsection*{Verification}

\begin{itemize}
\item \code{//lib/parts:parts\_test} drives two hosts and four:
\code{each\_answer\_goes\_home} has two hosts read at once from a
memory that answers with the address it was asked for, so a host given
another's beats would see an address it never asked for;
\code{write\_beats\_do\_not\_interleave} has both write four beats at
the same moment and checks the memory was given two whole bursts;
\code{four\_hosts\_share\_the\_link} has four read as fast as they can
for twelve hundred cycles and checks the counts come out within one of
each other.
\item \code{//lib/examples:ex\_axi\_arbiter} is the same shape as a
run, and the build simulates the netlist against it under nvc and
Verilator: \code{//docs:axi\_arbiter\_sim\_axi\_arbiter\_tb\_test} and
\code{//docs:axi\_arbiter\_vsim\_test}.
\end{itemize}

\subsection*{Used by}

\code{Board}, as \code{arb}, an \code{Arbiter4<32, 32, 4, 2, 4, 0>} in
front of its router: the core's tracker on port 0, the JTAG master's
pins on port 1 (issue 241), and the Ethernet port's fetch and store
engines on ports 2 and 3 (issue 151). Its peripheral side has four
bits of identifier, which four hosts fill, so a fifth widens the bus.

\subsection*{Limits}

The turn is three bits, which is what eight hosts need and no more.
\code{J} is not checked against \code{I} and the count. There is no
arbitration on quality of service, though the address phase carries
the field. A host that stops taking its answers stops the peripheral
side for every host, since the arbiter holds the beat rather than
dropping it.
